% Korean edition of figs/orderstates.tex.
\fig[wide]{orderstates}{서버가 받았을 수 있는 주문은 다시 보내지 않는다. 서버의 주문은 체결이든 정지든 \code{Done}으로 끝난다}[위: \code{AddOrder} 호출 하나가 응답 종류별로 하는 일. 아래: 서버에서 주문의 \code{state}(이름은 웹 클라이언트 코드에서)]{%
\begin{tikzpicture}[x=1mm,y=1mm,
  f/.style={font=\fontsize{7}{8.4}\selectfont},
  st/.style={rectangle,rounded corners=1.6mm,draw=navy,line width=0.6pt,f,text=navy,align=center,minimum width=24mm,minimum height=8.5mm,inner sep=1pt,fill=white},
  sv/.style={st,rounded corners=4.2mm,minimum width=22mm},
  un/.style={sv,draw=greyD,dashed,text=greyD},
  term/.style={st,fill=navy,text=white},
  tr/.style={-{Stealth[length=1.6mm]},line width=0.6pt,navy},
  trU/.style={-{Stealth[length=1.6mm]},line width=0.5pt,greyD,dashed},
  lb/.style={f,text=navy,inner sep=1pt,align=center},
  lbU/.style={f,text=greyD,inner sep=1pt,align=center},
  band/.style={font=\fontsize{7.5}{9}\selectfont\bfseries,text=navy,anchor=west}]
% ------------------------------------------------ wrapper band
\node[band] at (-22,33) {래퍼: \code{AddOrder} 호출 하나};
\node[st] (S) at (32,0) {\code{AddOrder}\\보내는 중};
\node[st] (F) at (84,0) {피드에서\\주문 찾는 중};
\node[st] (Q) at (134,0) {REST \code{Orders}\\조회 중};
\node[term] (T4) at (32,25) {\code{RequestError}\\발생};
\node[term] (TA) at (134,25) {\code{AmbiguousState}:\\주문 재전송 안 함};
\node[term] (TI) at (116,-17) {주문 ID 반환};
% self loops on S
\draw[tr] ([yshift=2mm]S.west) to[out=160,in=110,looseness=6] ([xshift=-7mm]S.north);
\node[lb,anchor=east,align=right] at (13,9) {401: 토큰 갱신 후\\한 번 더 보냄};
\draw[tr] ([yshift=-2mm]S.west) to[out=-160,in=-110,looseness=6] ([xshift=-7mm]S.south);
\node[lb,anchor=east,align=right] at (13,-9) {연결이 맺어지지 않음:\\한 번 더 보냄. 그것도\\실패하면 주문을 찾음};
% S -> T4, S -> F
\draw[tr] (S.north) -- node[lb,right]{400: 재시도 안 함} (T4.south);
\draw[tr] (S.east) -- node[lb,above]{주문 ID 없는\\200} node[lb,below,align=center]{또는 응답 유실:\\타임아웃, 5xx,\\보낸 뒤\\연결 끊김} (F.west);
% F -> Q, F -> TI, Q -> TA, Q -> TI
\draw[tr] (F.east) -- node[lb,above]{15\,s 안에} node[lb,below]{못 찾음} (Q.west);
\draw[tr] (Q.north) -- node[lb,right,align=left]{쓰기 이후 레코드에\\없음} (TA.south);
\draw[tr] (F.south) |- (TI.west);
\node[lb,anchor=north] at (92,-18.2) {찾음: 같은 종목코드,\\\code{ex}, 크기};
\draw[tr] (Q.south) |- node[lb,pos=0.25,right,xshift=0.8mm,align=left]{찾음} (TI.east);
% ------------------------------------------------ separator
\draw[greyM,line width=0.4pt] (-22,-27) -- (158,-27);
\node[band] at (-22,-31) {서버: 주문의 \code{state}};
% ------------------------------------------------ server band
\node[un] (G) at (-10,-44) {\code{Generated}};
\node[sv] (A) at (28,-44) {\code{Accepted}};
\node[sv] (W) at (70,-44) {\code{Working}};
\node[sv] (D) at (112,-44) {\code{Done}};
\node[un] (X) at (9,-62) {\code{Deleted}};
\draw[trU] (G) -- (A);
\draw[trU] (A) -- node[lbU,above,yshift=0.8mm]{관찰 못함} (W);
\draw[trU] (G) -- (X);
\draw[trU] (A) -- (X);
\node[lbU,anchor=east,align=right] at (-3,-58) {\code{DeleteOrder}};
\draw[tr] (W.north east) to[out=25,in=155] node[lb,above]{마지막 주식 체결} (D.north west);
\draw[tr] (W.east) -- node[lb,below,align=center]{정지} (D.west);
\draw[tr] (A.south) to[out=-35,in=-145] node[lb,below,align=center]{즉시 정지(주문 1650719)} (D.south);
\node[lb,anchor=west,align=left] at (125,-44) {\code{cnclT} null이면 체결,\\\code{cnclT} 설정이면 정지,\\거기에 \code{cumQty} $>$ 0이면\\일부 체결 후 정지};
% ------------------------------------------------ legend
\begin{scope}[shift={(-22,-78)}]
\node[st,minimum width=8mm,minimum height=4.5mm] at (4,0) {}; \node[f,anchor=west] at (9,0) {래퍼 단계};
\node[term,minimum width=8mm,minimum height=4.5mm] at (32,0) {}; \node[f,anchor=west] at (37,0) {호출 종료};
\node[sv,minimum width=8mm,minimum height=4.5mm,rounded corners=2.2mm] at (63,0) {}; \node[f,anchor=west] at (68,0) {서버 상태, 관찰함};
\node[un,minimum width=8mm,minimum height=4.5mm,rounded corners=2.2mm] at (101,0) {}; \node[f,anchor=west] at (106,0) {클라이언트 코드에만 있음};
\draw[tr] (145,1.6) -- (152,1.6); \node[f,anchor=west] at (153,1.6) {규칙 또는 관찰};
\draw[trU] (145,-1.8) -- (152,-1.8); \node[f,anchor=west] at (153,-1.8) {관찰 못함};
\end{scope}
\end{tikzpicture}}{래퍼는 첫 주문이 서버에 닿았는지 알 수 없으면 주문을 다시 보내지 않는다. 한 번 더 보내는 경우는 중복을 만들 수 없는 두 가지뿐이다. 연결이 아예 맺어지지 않은 경우(\code{ConnectError}, \code{ConnectTimeout}, \code{PoolTimeout})와, ASP.NET이 요청이 핸들러에 닿기 전에 돌려주는 401이다(D36). 연결 실패 뒤 재전송도 실패하면 주문은 피드·REST 검색으로 넘어가며, 세 번째 전송은 없다. 어떤 요청에서든 403이나 429를 받으면 \code{BlockedError}를 던지고 모든 쓰기를 멈춘다(그림에 없음). 잠금 하나가 보내기와 찾기를 덮고 이미 돌려준 ID는 건너뛰므로, 같은 주문 2건이 같은 주문 ID에 대응될 수 없다. 서버 상태: \code{Accepted}는 주문 1650719를 넣은 직후, \code{Working}은 08:30 주문에서 보았다. \code{Generated}와 \code{Deleted}는 클라이언트 코드에만 나온다. 끝난 주문은 모두 \code{Done}으로 보이므로 \code{cnclT}와 \code{cumQty}로 체결과 정지를 구분한다.}{\code{pytimefolio}의 \code{orders}, \code{transport} 모듈; 2026년 11월 24일 실시간 피드와 \code{Orders} 응답; DART Quant Team}
